\section{Proceso de entrenamiento e inferencia}

\subsection{Algoritmo de entrenamiento}

El algoritmo completo de entrenamiento de Conditional Flow Matching es el siguiente:

\begin{lstlisting}[caption={Algoritmo de entrenamiento de Conditional Flow Matching (ruta OT)}]
# v_theta: red neuronal que parametriza el campo vectorial
# p_data: distribución de datos (conjunto de entrenamiento)

for cada paso de entrenamiento:
    # 1. Muestrear datos y ruido
    x1 = sample_from(p_data)           # muestra de datos
    x0 = sample_from(Normal(0, I))     # muestra de ruido
    t  = sample_from(Uniform(0, 1))    # paso temporal

    # 2. Calcular interpolación (mapa de flujo condicional)
    xt = (1 - t) * x0 + t * x1

    # 3. Calcular objetivo (campo vectorial condicional)
    target = x1 - x0

    # 4. Calcular pérdida y actualizar
    prediction = v_theta(t, xt)
    loss = MSE(prediction, target)
    loss.backward()
    optimizer.step()
\end{lstlisting}

\begin{warningbox}{Precauciones durante el entrenamiento}
\begin{enumerate}
    \item El tiempo $t$ debe muestrearse uniformemente de $[0, 1]$; algunas implementaciones realizan un recorte pequeño en los extremos (por ejemplo, $[\epsilon, 1-\epsilon]$) para evitar problemas numéricos.
    \item Para la ruta OT, el objetivo $x_1 - x_0$ no depende de $t$, pero $x_t$ sí depende de $t$; por lo tanto, la red aún debe tener $t$ como entrada.
    \item La arquitectura de la red puede reutilizar la de U-Net o DiT de los modelos difusos, siempre y cuando se garantice que la codificación del tiempo $t$ sea correcta.
\end{enumerate}
\end{warningbox}

\subsection{Algoritmo de inferencia (muestreo)}

Después de completar el entrenamiento, el proceso para generar nuevas muestras es el siguiente:

\begin{lstlisting}[caption={Algoritmo de muestreo de Flow Matching (método Euler)}]
# v_theta: red neuronal del campo vectorial entrenada
# N: número de pasos de discretización
# dt: tamaño del paso = 1/N

# 1. Muestrear de la distribución de referencia
x = sample_from(Normal(0, I))

# 2. Resolver la ODE desde t=0 hasta t=1
for i in range(N):
    t = i / N
    x = x + v_theta(t, x) * dt    # paso de Euler

# x es ahora aproximadamente una muestra de p_data
return x
\end{lstlisting}

Se pueden utilizar solucionadores de ODE de orden superior (como Heun o Dormand-Prince/RK45) para mejorar la calidad del muestreo o reducir el número de pasos.

\subsection{Resumen de la sección}

Los procesos de entrenamiento e inferencia de Flow Matching son extremadamente concisos: el entrenamiento consiste en regresión de objetivos de velocidad en puntos de interpolación aleatorios, y la inferencia implica resolver una ODE partiendo del ruido. Todo el proceso no involucra ninguna simulación de SDE ni diseño complejo de programación de ruido.
